% diagrams/firstwords/funcion.tex -- la función de Navidad (el viernes que
% cierra el otoño): Lucía y Amy cantando "Jingle Bells" en el escenario,
% entre las cortinas, con las notas volando y dos cascabeles colgados.
\begin{tikzpicture}[dibujofw]
  % el escenario
  \filldraw[fill=white] (-5.2,-0.8) rectangle (10.2,0);
  \filldraw[fill=white] (-5.6,10.2) rectangle (10.6,11.0);
  % las cortinas
  \filldraw[fill=white] (-5.6,10.2) -- (-5.6,0) -- (-4.0,0) .. controls (-3.2,3.0) and (-2.4,8.0) .. (-2.8,10.2) -- cycle;
  \draw (-4.8,10.2) -- (-4.8,0.3); \draw (-3.9,10.2) .. controls (-3.8,6.0) .. (-3.3,1.2);
  \filldraw[fill=white] (10.6,10.2) -- (10.6,0) -- (9.0,0) .. controls (8.2,3.0) and (7.4,8.0) .. (7.8,10.2) -- cycle;
  \draw (9.8,10.2) -- (9.8,0.3); \draw (8.9,10.2) .. controls (8.8,6.0) .. (8.3,1.2);
  % Lucía y Amy, cantando (la boca en o)
  \foreach \x/\quien in {0/Lucia, 5.0/Amy} {
    \begin{scope}[shift={(\x,0)}]
      \ninaFW{\quien}{Abajo}{Abajo}
      \begin{scope}[shift={(0,6.3)}]
        \fill[white] (0,-0.58) ellipse (0.5 and 0.2);
        \filldraw[fill=white] (0,-0.55) ellipse (0.2 and 0.24);
      \end{scope}
    \end{scope}}
  % las notas y los cascabeles
  \foreach \x/\y/\a in {-1.8/8.8/0, 1.9/9.2/10, 3.2/8.0/-10, 6.9/8.9/5}
    {\begin{scope}[shift={(\x,\y)}, rotate=\a] \notaFW \end{scope}}
  \foreach \x in {-0.5,5.5} {\draw (\x,10.2) -- (\x,9.6); \begin{scope}[shift={(\x,9.0)}, scale=0.9] \cascabelFW \end{scope}}
\end{tikzpicture}
